%% Dimensionality from Self-Reference: Formalized Proof
%% To be integrated into spectral cosmology paper

\section{The Dimensional Constraint}

\subsection{Motivation}

The factor of 3 in $\Lambda = 3H^2$ appears as a derived consequence, not an input. We prove that $n = 3$ spatial dimensions is the unique dimension compatible with stable, grounded self-reference.

\subsection{Definitions}

\begin{definition}[Integration Capacity]
For a region $S$ of characteristic scale $L$ in a relational structure with Laplacian $\mathcal{L}$, the \emph{integration capacity} is:
\[
\mathcal{I}(S) = \frac{1}{\lambda_1(\mathcal{L}|_S)}
\]
By Weyl's law, $\lambda_1 \geq c_n / L^2$ for regions of size $L$ in $n$ dimensions, giving:
\[
\mathcal{I}(S) \leq \frac{L^2}{c_n}
\]
\end{definition}

\begin{definition}[Boundary Complexity]
For a region $S \subset M$ in an $n$-dimensional manifold, the \emph{boundary complexity} is:
\[
\mathcal{B}(S) = |\partial S| \sim L^{n-1}
\]
\end{definition}

\begin{definition}[Grounded Self-Reference]
A subsystem $S$ exhibits \emph{grounded self-reference} if it can:
\begin{enumerate}
    \item Model its own internal dynamics (integration)
    \item Distinguish itself from its environment (boundary representation)
    \item Maintain this distinction under perturbation (stability)
\end{enumerate}
\end{definition}

\subsection{The Main Theorem}

\begin{theorem}[Dimensionality from Self-Reference] \label{thm:dimension}
\textsc{[Derived]}

A relational structure admits stable grounded self-referential subsystems at all scales if and only if the spatial dimension $n \leq 3$. The case $n = 3$ uniquely saturates the bound, maximizing complexity compatible with self-reference.
\end{theorem}

\begin{proof}
Let $S$ be a self-referential subsystem of characteristic scale $L$.

\textbf{Step 1: Integration requirement.}
Self-modeling requires integrating information across $S$. The time scale for integration is $\tau \sim L^2 / D$ (diffusion). The capacity for this integration is bounded by the spectral gap:
\[
\mathcal{I}(S) \leq \frac{L^2}{c_n}
\]

\textbf{Step 2: Boundary requirement.}
Grounded self-reference requires modeling the interface $\partial S$ between self and environment. The information required scales with boundary size:
\[
\mathcal{B}(S) \sim L^{n-1}
\]

\textbf{Step 3: Grounding condition.}
For $S$ to model itself as embedded in its environment:
\[
\mathcal{I}(S) \geq \mathcal{B}(S)
\]
\[
L^2 \geq c \cdot L^{n-1}
\]

\textbf{Step 4: Scale independence.}
For this to hold at all scales $L$ (not just small $L$):
\[
2 \geq n - 1 \implies n \leq 3
\]

\textbf{Step 5: Saturation analysis.}
\begin{itemize}
    \item $n = 1$: No interior. Boundary is 0-dimensional (two points). Self/other distinction trivial. Insufficient structure.
    \item $n = 2$: Interior exists but $L^2 \gg L^1$ for large $L$. System over-resourced. Self-modeling ``too easy'' — no optimization pressure.
    \item $n = 3$: $L^2 \sim L^2$. Integration capacity exactly matches boundary complexity. Every integration degree of freedom corresponds to a boundary degree of freedom. Maximum richness.
\end{itemize}

\textbf{Step 6: Holographic consistency.}
The Bekenstein bound gives $S_{max} \sim L^{n-1}$. If self-modeling requires representing all degrees of freedom defining the system, this independently requires capacity $\geq L^{n-1}$, confirming $n \leq 3$.
\end{proof}

\begin{corollary}[The Cosmological Coefficient] \label{cor:coefficient}
\textsc{[Derived]}

In $n = 3$ spatial dimensions, the cosmological constant satisfies $\Lambda = 3H^2$, where the factor of 3 is the spatial dimension.
\end{corollary}

\begin{proof}
For de Sitter space with cosmological constant $\Lambda$ in $(n+1)$ spacetime dimensions:
\begin{itemize}
    \item Spatial section is $S^n$ with radius $a^2 = n(n-1)/(2\Lambda)$
    \item First eigenvalue: $\lambda_1 = n/a^2 = 2\Lambda/(n-1)$
    \item For $n = 3$: $\lambda_1 = 2\Lambda/2 = \Lambda$
    \item Hubble parameter: $H^2 = \Lambda/3$
    \item Therefore: $\lambda_1 = \Lambda = 3H^2$
\end{itemize}
\end{proof}

\subsection{Connection to Self-Alignment}

\begin{proposition}[Dimensional Selection Principle]
\textsc{[Conjectured]}

Among possible spatial dimensions, $n = 3$ maximizes the cosmic self-alignment score $S_{cosmic}$.
\end{proposition}

\begin{remark}
This provides a variational principle for dimensional selection: the universe realizes the dimension that optimizes self-reference. Lower dimensions have excess integration capacity relative to boundary complexity, resulting in lower self-alignment scores (less efficient use of resources). Higher dimensions cannot sustain stable self-reference.
\end{remark}
